% ============================================================
%  APPENDIX A --- Installation
% ============================================================
\chapter{Instalación}
\label{app:installation}

Este apéndice cubre todos los escenarios de implementación compatibles para FairWindSK,
de una fuente de primera vez construir en un ordenador portátil desarrollador a un completamente automático
bota de quiosco en una pantalla de helm dedicada.
Cada escenario lista requisitos de hardware o plataforma, comandos exactos,
y pasos de verificación de post-instalación.

\begin{notebox}
Todos los escenarios suponen un servidor Signal~K en algún lugar del barco
Red Wi-Fi.
FairWindSK es un shell de pantalla, no un servidor de datos.
Instala y configura Signal~K antes de instalar FairWindSK.
\end{notebox}

% ── A.1  Prerequisites ───────────────────────────────────────
\section{Prerrequisitos}
\label{app:prerequisites}

\subsection{Requisitos comunes para todas las plataformas de escritorio}

\begin{itemize}
  \item CMake 3.16 or newer
  \item A C++17-capable compiler supported by the chosen Qt~6 kit
  \item Git
  \item Internet access during the first build
        (CMake downloads \code{nlohmann/json}, \code{QtZeroConf}, and \code{QHotkey}
automáticamente; las obras posteriores utilizan las copias en caché)
\end{itemize}

\subsection{Se requiere Módulos Qt 6 --- escritorio}

\rowcolors{2}{LtGray}{white}
\begin{longtable}{L{4.5cm}L{9.2cm}}
\toprule
\rowcolor{Navy}
\color{white}\sffamily\bfseries Module &
\color{white}\sffamily\bfseries Purpose \\
\midrule
\code{Core}            & Qt base classes and event loop. \\
\code{Gui}             & 2D graphics and font handling. \\
\code{Widgets}         & Widget toolkit used throughout the shell. \\
\code{Concurrent}      & Thread pool for background tasks. \\
\code{Network}         & HTTP/REST client and TLS. \\
\code{WebSockets}      & Signal~K WebSocket stream. \\
\code{Xml}             & Configuration and resource parsing. \\
\code{Svg}             & OpenBridge icon rendering. \\
\code{Positioning}     & GPS location APIs (desktop fallback). \\
\code{WebEngineWidgets}& Embedded browser for hosted Signal~K web apps. \\
\code{VirtualKeyboard} & On-screen keyboard for touchscreen deployments. \\
\code{PrintSupport}    & PDF generation (crash reports). \\
\bottomrule
\end{longtable}

\subsection{Se requiere Módulos Qt 6 --- Android e iOS}

Mobile builds replace \code{WebEngineWidgets} with a lighter web layer and
dejar las dependencias de escritorio solamente:

\rowcolors{2}{LtGray}{white}
\begin{longtable}{L{4.5cm}L{9.2cm}}
\toprule
\rowcolor{Navy}
\color{white}\sffamily\bfseries Module &
\color{white}\sffamily\bfseries Replaces or adds \\
\midrule
\code{Qml}         & QML runtime for WebView hosting. \\
\code{Quick}       & Qt Quick scene graph. \\
\code{QuickWidgets}& \code{QQuickWidget} bridge between widget and QML worlds. \\
\code{WebView}     & Replaces \code{WebEngineWidgets}. Uses the platform's native browser engine. \\
\bottomrule
\end{longtable}

Mobile builds do \emph{not} include \code{QtZeroConf}, \code{QHotkey},
\code{PrintSupport}, or \code{WebEngineWidgets}.

% ── A.2  Generic desktop build ───────────────────────────────
\section{Generic Desktop Build (todas las plataformas)}
\label{app:generic}

Este es el procedimiento de referencia que se aplica a cada plataforma de escritorio
antes de añadir pasos específicos de la plataforma.

\begin{enumerate}
  \item Clone the repository:
\begin{lstlisting}
git clone https://github.com/OpenFairWind/FairWindSK.git
cd FairWindSK
\end{lstlisting}

  \item Configure:
\begin{lstlisting}
cmake -S . -B build
\end{lstlisting}

  \item Build (parallel across all CPU cores):
\begin{lstlisting}
cmake --build build --parallel
\end{lstlisting}

  \item Install to the default prefix:
\begin{lstlisting}
cmake --install build
\end{lstlisting}

  \item Override the install prefix if needed:
\begin{lstlisting}
cmake -S . -B build -DCMAKE_INSTALL_PREFIX=/desired/prefix
cmake --build build --parallel
cmake --install build
\end{lstlisting}

  \item If Qt is installed outside the default search path, pass
        \code{-DCMAKE\_PREFIX\_PATH} during configure:
\begin{lstlisting}
cmake -S . -B build -DCMAKE_PREFIX_PATH=/path/to/Qt/6.x.x/gcc_64
\end{lstlisting}
\end{enumerate}

\begin{tipbox}
Si el acceso a Internet está restringido durante la construcción, instalar
\code{nlohmann\_json} from your package manager first.
CMake utilizará el encabezado del sistema en lugar de descargarlo.
On Debian/Ubuntu: \code{sudo apt install nlohmann-json3-dev}.
On macOS with Homebrew: \code{brew install nlohmann-json}.
\end{tipbox}

% ── A.3  macOS ───────────────────────────────────────────────
\section{macOS Instalación}
\label{app:macos}

\subsection{Necesidades}

\begin{itemize}
  \item macOS 12 Monterey or newer (recommended)
  \item Xcode command line tools: \code{xcode-select --install}
  \item Qt~6 desktop kit (via Qt Online Installer or Homebrew)
  \item CMake and Ninja (via Homebrew or Qt Maintenance Tool)
\end{itemize}

\subsection{Instalar con Homebrew}

\begin{lstlisting}
brew install qt cmake ninja nlohmann-json
cmake -S . -B build -G Ninja \
  -DCMAKE_PREFIX_PATH="$(brew --prefix qt)"
cmake --build build --parallel
\end{lstlisting}

\subsection{Corre desde el árbol de la construcción}

\begin{lstlisting}
open build/FairWindSK.app
\end{lstlisting}

\subsection{Instala al sistema}

\begin{lstlisting}
cmake --install build
# Installs to /usr/local by default:
open /usr/local/FairWindSK.app
\end{lstlisting}

\subsection{Redistributo (distribución) Qt runtime)}

Para distribución fuera del árbol de construcción, utilice la herramienta de implementación de Apple:

\begin{lstlisting}
macdeployqt build/FairWindSK.app
\end{lstlisting}

This bundles all Qt frameworks inside the \code{.app} so the app runs
en máquinas sin Qt instalado.

\subsection{Verificación posterior a la instalación}

\begin{enumerate}
  \item Launch FairWindSK.
  \item Open \ui{Settings~> Connection}.
        Enter your Signal~K server address and tap \key{Connect}.
  \item Confirm the status changes to \textit{Live} and apps appear on the launcher.
\end{enumerate}

% ── A.4  Linux desktop ───────────────────────────────────────
\section{Instalación de escritorio}
\label{app:linux}

\subsection{Requisitos (Debian / Ubuntu)}

\begin{lstlisting}
sudo apt update
sudo apt install \
  build-essential cmake ninja-build git \
  nlohmann-json3-dev \
  qt6-base-dev qt6-base-dev-tools \
  qt6-webengine-dev qt6-webengine-dev-tools \
  qt6-websockets-dev qt6-positioning-dev \
  libqt6svg6-dev qt6-virtualkeyboard-dev \
  qt6-base-private-dev libqt6printsupport6 \
  libavahi-compat-libdnssd-dev libnss-mdns avahi-utils
\end{lstlisting}

\subsection{Construir y correr}

\begin{lstlisting}
cmake -S . -B build -G Ninja
cmake --build build --parallel
./build/FairWindSK
\end{lstlisting}

\subsection{Sistema de instalación}

\begin{lstlisting}
sudo cmake --install build
\end{lstlisting}

\code{cmake --install} on a Linux desktop target:

\begin{itemize}
  \item Installs the \code{FairWindSK} binary to
        \code{\$\{CMAKE\_INSTALL\_BINDIR\}} (default \code{/usr/local/bin}).
  \item Installs the bundled icon directory alongside the binary.
  \item Installs the desktop-only helper libraries to
        \code{\$\{CMAKE\_INSTALL\_LIBDIR\}}.
  \item Installs the \code{fairwindsk.desktop} launcher file to the standard
Directorio de aplicaciones XDG.
  \item Installs the \code{fairwindsk.png} menu icon to
        \code{/usr/local/share/icons/hicolor/256x256/apps/}.
\end{itemize}

Después de la instalación, FairWindSK aparece en el menú de aplicaciones.
Ejecutarlo desde el menú o con:

\begin{lstlisting}
FairWindSK
\end{lstlisting}

\subsection{Embalaje}

Incluir el tiempo de ejecución Qt y las dependencias externas construidas bajo
\code{build/external/lib} in your package.
The build-tree rpath means the binary runs from \code{build/} without
\code{LD\_LIBRARY\_PATH}, but a packaged install needs those libraries in
la ruta de la biblioteca del sistema o junto al binario.

% ── A.5  Raspberry Pi OS --- nav station ───────────────────────
\section{Raspberry Pi OS --- Estación de navegación}
\label{app:rpi-station}

Este escenario instala FairWindSK en un Raspberry~Pi usado como normal
estación de navegación de escritorio --- pantalla, teclado, y ratón adjunto,
dirigiendo el escritorio Raspberry~Pi~OS.

\subsection{Necesidades}

Raspberry Pi~4 o~5 con Raspberry~Pi~OS Bookworm (64-bit recomendado),
y una conexión estable a Internet para la descarga inicial del paquete.

\subsection{Instalar dependencias de construcción}

\begin{lstlisting}
sudo apt update
sudo apt install \
  build-essential cmake ninja-build git \
  nlohmann-json3-dev \
  qml6-module-qt-labs-folderlistmodel \
  qml6-module-qtquick-window \
  qml6-module-qtquick-layouts \
  qml6-module-qtqml-workerscript \
  libnss-mdns avahi-utils \
  libavahi-compat-libdnssd-dev libxkbcommon-dev \
  qt6-base-dev qt6-base-dev-tools \
  qt6-webengine-dev qt6-webengine-dev-tools \
  qt6-websockets-dev qt6-positioning-dev \
  libqt6svg6-dev \
  qt6-virtualkeyboard-dev \
  libqt6virtualkeyboard6 \
  qt6-virtualkeyboard-plugin
\end{lstlisting}

\begin{notebox}
Si la construcción falla con un desaparecido
\code{Qt6VirtualKeyboardConfigVersionImpl.cmake}, apply the known
todo el trabajo:
\begin{lstlisting}
sudo touch /usr/lib/aarch64-linux-gnu/cmake/Qt6VirtualKeyboard/\
Qt6VirtualKeyboardConfigVersionImpl.cmake
\end{lstlisting}
Then re-run \code{cmake}.
\end{notebox}

\subsection{Construir e instalar}

\begin{lstlisting}
git clone https://github.com/OpenFairWind/FairWindSK.git
cd FairWindSK
cmake -S . -B build -G Ninja
cmake --build build --parallel
sudo cmake --install build
\end{lstlisting}

FairWindSK aparece en el menú de aplicación bajo
\textit{Navigation} or \textit{Internet}.

\subsection{Integración OpenPlotter}

If OpenPlotter is detected during \code{cmake --install}, FairWindSK
También copia su icono de lanzamiento en los directorios del menú OpenPlotter.
No se requieren pasos adicionales.

% ── A.6  Raspberry Pi OS --- kiosk helm display ────────────────
\section{Raspberry Pi OS --- Kiosk Helm Display}
\label{app:rpi-kiosk}

Este es el escenario recomendado para una pantalla de helm dedicada:
a Raspberry~Pi~4 o~5 con pantalla táctil, ejecutando FairWindSK automáticamente
en el arranque sin teclado y sin menú de escritorio visible.

\subsection{Sinopsis}

\begin{enumerate}
  \item Install FairWindSK as in Section~\ref{app:rpi-station}.
  \item Configure FairWindSK for full-screen kiosk mode.
  \item Install the \code{fairwindsk-startup} helper to manage autostart.
  \item (Optional) Configure the desktop panel to auto-hide.
  \item Reboot and verify.
\end{enumerate}

\subsection{Paso 1 --- Construir e instalar}

Follow Section~\ref{app:rpi-station}.
On Raspberry~Pi~OS, \code{cmake --install} automatically installs the
\code{fairwindsk-startup.desktop} autostart entry to
\code{/etc/xdg/autostart/}.
FairWindSK comenzará automáticamente en el próximo login.

\subsection{Paso 2 --- Configurar el modo de quiosco de pantalla completa}

Open \ui{Settings~> Main} and set \key{Window Mode} to
\textbf{Full Screen}.
Esto requiere una ventana sin marco, siempre en la parte superior que cubre todo
pantalla sin barra de título o decoración de ventana.

Alternatively, edit \code{fairwindsk.json} directly:

\begin{lstlisting}
{
  "main": {
    "windowMode": "fullscreen"
  }
}
\end{lstlisting}

\subsection{Paso 3... Instalar el ayudante de arranque}

The \code{fairwindsk-startup} Python script manages the FairWindSK
proceso ciclo de vida.
It is installed by \code{cmake --install} alongside the binary.

What \code{fairwindsk-startup} does at each boot:

\begin{enumerate}
  \item Reads \code{fairwindsk.ini} to find the configuration path.
  \item Reads \code{fairwindsk.json} to check the \code{virtualKeyboard}
        flag and sets \code{QT\_IM\_MODULE=qtvirtualkeyboard} accordingly.
  \item Waits up to 45 seconds for the Signal~K web-app catalogue
        endpoint to respond (\code{/signalk/v1/apps/list} or the legacy
        \code{/skServer/webapps}).
Esto asegura que los iconos plugin y el arte están disponibles antes
FairWindSK comienza.
  \item Launches FairWindSK.
  \item If FairWindSK exits with code~1 (the \textit{Restart} button in
        \ui{Settings~> System}), the helper relaunches it immediately.
  \item If FairWindSK exits with code~0 (the \textit{Quit} button or
un cierre limpio), el ayudante sale limpiamente.
\end{enumerate}

\begin{warningbox}
Use \key{Restart FairWindSK} in \ui{Settings~> System} (exit code~1)
para relanzar después de cambios de configuración.
\key{Quit FairWindSK} (exit code~0) leaves the display blank
sin reiniciar.
\end{warningbox}

\subsection{Paso 4... Ocultar el panel de escritorio (labwc / Bookworm)}

En Raspberry~Pi~OS Bookworm con el compositor de labwc, la barra de tareas
panel puede sobreponer la ventana FairWindSK de pantalla completa.
Para evitarlo, configura el panel para auto-hide:

\begin{lstlisting}
# Edit the panel configuration file
nano ~/.config/wf-panel-pi.ini
\end{lstlisting}

Añadir o cambiar:

\begin{lstlisting}
[panel]
autohide = true
\end{lstlisting}

Guarda el archivo y reinicia.

\subsection{Paso 5 --- Activar teclado virtual (sólo pantalla táctil)}

Si la pantalla de helm no tiene teclado físico, active el Qt virtual
teclado:

\begin{enumerate}
  \item Open \ui{Settings~> Main}.
  \item Tick \key{Virtual Keyboard}.
  \item Tap \key{Restart FairWindSK} in \ui{Settings~> System}.
\end{enumerate}

El ayudante de arranque lee el ajuste en cada bota y establece el
variable ambiente automáticamente.

\subsection{Paso 6 --- Verificar la bota del quiosco}

\begin{enumerate}
  \item Reboot the Raspberry~Pi.
  \item Verify FairWindSK starts automatically (allow up to 60 seconds
Para Signal~K y FairWindSK empezar juntos).
  \item Confirm the display is full-screen with no title bar.
  \item Confirm the status shows \textit{Live}.
  \item Confirm apps appear on the launcher.
\end{enumerate}

\begin{tipbox}
En un Raspberry~Pi~4 con una tarjeta microSD clase~A2, FairWindSK
normalmente se vuelve interactivo en menos de 30 segundos de power-on,
suponiendo que Signal~K comience rápidamente.
Use a fast card (\textgreater{}90~MB/s read) to minimise boot time.
\end{tipbox}

% ── A.7  Windows ─────────────────────────────────────────────
\section{Instalación de Windows}
\label{app:windows}

\subsection{Necesidades}

\begin{itemize}
  \item Windows~10 or~11 (64-bit)
  \item Qt~6 desktop kit for MSVC
(instalado a través de Qt Online Installer)
  \item Visual Studio Build Tools 2022 or Visual Studio Community 2022
  \item CMake and Ninja
(con Qt o instalados por separado)
\end{itemize}

\subsection{Construir e instalar}

Abra un programa de desarrollo con el entorno MSVC cargado,
entonces:

\begin{lstlisting}
git clone https://github.com/OpenFairWind/FairWindSK.git
cd FairWindSK
cmake -S . -B build -G Ninja ^
  -DCMAKE_PREFIX_PATH="C:\Qt\6.x.x\msvc2022_64"
cmake --build build --parallel
cmake --install build
\end{lstlisting}

\subsection{Despliegue el tiempo de ejecución Qt}

Después de una exitosa construcción, desplegar los DLLs Qt runtime junto a
ejecutable:

\begin{lstlisting}
windeployqt build\FairWindSK.exe
\end{lstlisting}

\code{cmake --install} installs the executable, the bundled icon
directory under \code{bin\textbackslash icons}, and the desktop-only
helper DLLs into the same \code{bin} directory.

\subsection{Corre}

\begin{lstlisting}
build\FairWindSK.exe
\end{lstlisting}

\begin{notebox}
\code{file://} launcher entries are retained in configuration for
compatibilidad pero están bloqueados por el modelo de ventanilla única en Windows.
Use \code{http://} or \code{https://} URLs for all applications.
\end{notebox}

% ── A.8  Android ─────────────────────────────────────────────
\section{Instalación de Android}
\label{app:android}

Android construye utilizar la implementación móvil Qt WebView.
Qt Creator es el flujo de trabajo recomendado para Android porque administra
el SDK, NDK, la firma de APK y el despliegue de dispositivos.

\subsection{Necesidades}

\begin{itemize}
  \item A macOS, Linux, or Windows build host
  \item Qt~6 with the Android kit for your target ABI
(arm64-v8a se recomienda para dispositivos modernos)
  \item Android Studio (for the SDK, NDK, and emulator)
  \item Java JDK (bundled with Android Studio or installed separately)
  \item CMake and Ninja
\end{itemize}

\subsection{Guía de construcción paso a paso}

\begin{enumerate}
  \item Open Android Studio and install:
una plataforma SDK Android, SDK Platform-Tools, SDK Build-Tools,
y la versión NDK recomendada por su versión Qt.
  \item In the Qt Maintenance Tool, install the Android Qt components
que coincide con su anfitrión y objetivo ABI.
  \item In Qt Creator, go to
        \ui{Preferences~> Devices~> Android} and verify that the SDK,
NDK, JDK y las vías de la plataforma se detectan correctamente.
  \item Clone the repository:
\begin{lstlisting}
git clone https://github.com/OpenFairWind/FairWindSK.git
\end{lstlisting}
  \item Open \code{FairWindSK/CMakeLists.txt} in Qt Creator.
  \item Select an Android kit (e.g.\ \textit{Android Qt~6.x.x Clang arm64-v8a}).
  \item Let Qt Creator configure the project.
        The mobile configuration uses \code{Qt::WebView},
        \code{Qt::Quick}, and \code{Qt::QuickWidgets},
y salta las dependencias de escritorio solo.
  \item Build with \ui{Build~> Build Project}.
  \item Connect an Android device with developer mode enabled,
o iniciar un emulador Android.
  \item Deploy with \ui{Build~> Run}.
\end{enumerate}

\subsection{Configuración de línea de comando (avanzado)}

\begin{lstlisting}
cmake -S . -B build-android -G Ninja \
  -DCMAKE_TOOLCHAIN_FILE="$ANDROID_NDK/build/cmake/android.toolchain.cmake" \
  -DANDROID_ABI=arm64-v8a \
  -DANDROID_PLATFORM=android-24 \
  -DCMAKE_PREFIX_PATH="/path/to/Qt/6.x.x/android_arm64_v8a"
cmake --build build-android --parallel
\end{lstlisting}

Qt Creator sigue siendo recomendado para el embalaje y el despliegue de dispositivos.

\subsection{Android características diferencias}

\rowcolors{2}{LtGray}{white}
\begin{longtable}{L{5.2cm}L{8.5cm}}
\toprule
\rowcolor{Navy}
\color{white}\sffamily\bfseries Feature &
\color{white}\sffamily\bfseries Android status \\
\midrule
Hosted web apps          & Supported via Qt WebView. \\
Signal~K stream          & Supported (WebSockets). \\
mDNS server discovery    & Not available (disabled on Android). \\
Global hotkey SHIFT+TAB  & Not available (desktop-only, uses QHotkey). \\
Virtual keyboard         & Uses Android system keyboard; Qt Virtual Keyboard setting has no effect. \\
File printing / PDF      & Not available (PrintSupport not included). \\
QWebEngine cookies       & Not used; mobile path uses a separate mechanism. \\
\code{file://} apps      & Blocked by the single-window model on all targets. \\
\bottomrule
\end{longtable}

% ── A.9  iOS and iPadOS ──────────────────────────────────────
\section{Instalación iOS y iPadOS}
\label{app:ios}

iOS y iPadOS construyen utilizar la misma aplicación web móvil como Android.
Qt Creator con una cuenta de desarrollador de Apple es necesario para el despliegue de dispositivos.

\subsection{Necesidades}

\begin{itemize}
  \item A macOS build host
  \item Xcode and Xcode command line tools
  \item Qt~6 with the iOS kit
  \item CMake and Ninja
  \item An Apple developer account and provisioning profile
(para el despliegue de dispositivos; no es necesario para el Simulador)
\end{itemize}

\subsection{Guía de construcción paso a paso}

\begin{enumerate}
  \item Open Xcode at least once and accept the license agreement.
  \item In the Qt Maintenance Tool, install the iOS Qt components
que coincide con el escritorio Qt liberación.
  \item Clone the repository:
\begin{lstlisting}
git clone https://github.com/OpenFairWind/FairWindSK.git
\end{lstlisting}
  \item Open \code{FairWindSK/CMakeLists.txt} in Qt Creator.
  \item Select an iOS kit (e.g.\ \textit{Qt~6.x.x for iOS}).
  \item Build with \ui{Build~> Build Project}.
  \item To run on the Simulator: choose an iOS Simulator run target
        and select \ui{Build~> Run}.
  \item To run on a device: connect the iPhone or iPad, select the
dispositivo físico en el selector de destino de ejecución, verificar la configuración de firma
cuando fue incitado, y corrió desde Qt Creator.
\end{enumerate}

\subsection{Configuración de línea de comando (avanzado)}

\begin{lstlisting}
cmake -S . -B build-ios -G Ninja \
  -DCMAKE_SYSTEM_NAME=iOS \
  -DCMAKE_OSX_SYSROOT=iphoneos \
  -DCMAKE_OSX_ARCHITECTURES=arm64 \
  -DCMAKE_PREFIX_PATH="/path/to/Qt/6.x.x/ios"
cmake --build build-ios --parallel
\end{lstlisting}

Use \code{iphonesimulator} and \code{x86\_64} (or \code{arm64} on
Apple Silicon hosts) para Simulator construye.

\subsection{Diferencias de características iOS}

\rowcolors{2}{LtGray}{white}
\begin{longtable}{L{5.2cm}L{8.5cm}}
\toprule
\rowcolor{Navy}
\color{white}\sffamily\bfseries Feature &
\color{white}\sffamily\bfseries iOS / iPadOS status \\
\midrule
Hosted web apps          & Supported via Qt WebView (WKWebView). \\
Signal~K stream          & Supported (WebSockets). \\
mDNS server discovery    & Not available. Enter the server URL manually. \\
Global hotkey SHIFT+TAB  & Not available. \\
Virtual keyboard         & Uses iOS system keyboard. \\
File printing / PDF      & Not available. \\
\code{file://} apps      & Blocked by the single-window model on all targets. \\
\bottomrule
\end{longtable}

% ── A.10  Signal K Server on Docker ──────────────────────────
\section{Signal K Server en Docker}
\label{app:signalk-docker}

Si ya no tienes un servidor Signal~K, Docker es el más rápido
forma de conseguir uno corriendo.
Esta sección documenta el despliegue de Docker desde el
Repositorio del proyecto FairWindSK.

\subsection{Docker image}

\begin{lstlisting}
cr.signalk.io/signalk/signalk-server:latest
\end{lstlisting}

Supported architectures: \code{linux/amd64}, \code{linux/arm/v7},
\code{linux/arm64} (covers all Raspberry~Pi~4 and~5 configurations).

\subsection{Inicio rápido con datos demo NMEA}

\begin{lstlisting}
mkdir $HOME/.signalk
docker run --init -it --rm \
  --name signalk-server \
  --publish 3000:3000 \
  -v $HOME/.signalk:/home/node/.signalk \
  cr.signalk.io/signalk/signalk-server \
  --sample-nmea0183-data
\end{lstlisting}

Esto ejecuta Signal~K en port~3000 con datos simulados NMEA~0183.
Úsalo para probar FairWindSK sin instrumentos reales.

\subsection{Despliegue de producción}

\begin{lstlisting}
docker run -d --init \
  --name signalk-server \
  -p 3000:3000 \
  -v $HOME/.signalk:/home/node/.signalk \
  cr.signalk.io/signalk/signalk-server
\end{lstlisting}

\begin{itemize}
  \item The \code{-d} flag runs the container as a background daemon.
  \item The \code{-v} flag mounts the configuration directory so
Ajustes Signal~K, plugins y datos persisten a través de contenedores
reinicia.
  \item On first access, the Signal~K admin web UI at
        \code{http://localhost:3000} prompts you to create admin credentials.
\end{itemize}

\subsection{Actualización de la señal K}

\begin{lstlisting}
docker pull cr.signalk.io/signalk/signalk-server:latest
docker stop signalk-server
docker rm signalk-server
# Re-run the production docker run command above
\end{lstlisting}

\subsection{Estructura del directorio dentro del contenedor}

\rowcolors{2}{LtGray}{white}
\begin{longtable}{L{4.8cm}L{8.9cm}}
\toprule
\rowcolor{Navy}
\color{white}\sffamily\bfseries Path &
\color{white}\sffamily\bfseries Contents \\
\midrule
\code{/home/node/signalk}  & Signal~K server files (read-only inside container). \\
\code{/home/node/.signalk} & Settings files, plugins, and data.
                              \textbf{Mount this to the host.} \\
\bottomrule
\end{longtable}

% ── A.11  First-Run Configuration ────────────────────────────
\section{Configuración de primer orden}
\label{app:firstrun}

Después de la instalación en cualquier plataforma, siga estos pasos en el primer lanzamiento.

\subsection{Localizaciones de archivos de configuración}

FairWindSK almacena su configuración en dos archivos en el usuario
directorio de configuración Qt:

\rowcolors{2}{LtGray}{white}
\begin{longtable}{L{3.6cm}L{10.1cm}}
\toprule
\rowcolor{Navy}
\color{white}\sffamily\bfseries File &
\color{white}\sffamily\bfseries Contents \\
\midrule
\code{fairwindsk.json}  &
Todos los ajustes visuales del usuario: URL de conexión, diseños de barras, catálogo de aplicaciones,
definiciones de widget, presets de confort, unidades y mapas de ruta Signal~K.
JSON legible por humanos.
  Safe to back up, transfer, and version-control. \\
\code{fairwindsk.ini}   &
  Qt-managed settings: path to \code{fairwindsk.json}, the debug flag,
y la autenticación token.
El token se mantiene aquí intencionadamente así que nunca se incluye en un
  configuration export. \\
\bottomrule
\end{longtable}

On first run, if \code{fairwindsk.json} does not exist, FairWindSK
semillas de los predeterminados agrupados en
\code{resources/json/configuration.json}.

\subsection{Conectar a la señal K}

\begin{enumerate}
  \item Open \ui{Settings~> Connection}.
  \item Enter the Signal~K server address,
        e.g.\ \code{http://192.168.1.100:3000}.
Si estás ejecutando Signal~K en el mismo dispositivo,
        use \code{http://localhost:3000}.
  \item Tap \key{Connect}.
  \item Wait for the status to change to \textit{Live}.
  \item If the server requires authentication,
        tap \key{Request Token} and approve the request in the
Interfaz de administración Signal~K.
\end{enumerate}

\subsection{Authenticate (opcional)}

Si el servidor Signal~K tiene la seguridad activada:

\begin{enumerate}
  \item Tap \key{Request Token} in \ui{Settings~> Connection}.
  \item FairWindSK opens the Signal~K admin interface in the
Bottom~Bar cajón, navegado a
        \ui{Security~> Access Requests}.
  \item Approve the request as an admin user.
  \item FairWindSK stores the token automatically and reconnects.
\end{enumerate}

\subsection{Ajustes esenciales de primera hora}

\rowcolors{2}{LtGray}{white}
\begin{longtable}{L{4.0cm}L{9.7cm}}
\toprule
\rowcolor{Navy}
\color{white}\sffamily\bfseries Setting &
\color{white}\sffamily\bfseries Where and what to set \\
\midrule
Unidades&
  \ui{Settings~> Units}.
  Choose knots (speed), nautical miles (distance), and metres or feet (depth). \\
Formato de coordinación&
  \ui{Settings~> Main}.
  Choose DD°MM.MMM' for standard nautical use. \\
Modo de ventana&
  \ui{Settings~> Main}.
Elige Full~Screen para una pantalla de helm dedicada;
  Windowed or Centred for a navigation station. \\
Confort preset&
  \ui{Settings~> Comfort}.
  Set Day for initial testing; change to Night for the night watch. \\
Aplicaciones de lanzamiento&
  \ui{Settings~> Applications}.
  Arrange your chart plotter, instruments, and other apps on page~1. \\
teclado virtual&
  \ui{Settings~> Main} (touchscreen only).
  Tick, then restart FairWindSK. \\
\bottomrule
\end{longtable}

% ── A.12  Configuration file reference ───────────────────────
\section{Campos de configuración clave}
\label{app:config}

Advanced users and system integrators can edit \code{fairwindsk.json}
directamente.
Esta sección abarca los campos más útiles.

\subsection{Conexión}

\begin{lstlisting}
{
  "connection": {
    "server": "http://your-signalk-host:3000"
  }
}
\end{lstlisting}

\subsection{Modo de ventana}

\begin{lstlisting}
{
  "main": {
    "windowMode": "fullscreen",
    "language": "en",
    "virtualKeyboard": false
  }
}
\end{lstlisting}

Valid values for \code{windowMode}: \code{windowed}, \code{centered},
\code{maximized}, \code{fullscreen}.
Valid values for \code{language}: \code{system}, \code{en}, \code{fr},
\code{es}, \code{it}.

\subsection{Unidades}

\begin{lstlisting}
{
  "units": {
    "vesselSpeed": "kn",
    "windSpeed":   "kn",
    "distance":    "nm",
    "depth":       "mt"
  }
}
\end{lstlisting}

\subsection{La trayectoria de Signal K se anula}

\begin{lstlisting}
{
  "signalk": {
    "sog": "navigation.speedOverGround",
    "dpt": "environment.depth.belowTransducer",
    "notifications.pob": "notifications.mob"
  }
}
\end{lstlisting}

\subsection{Diagnosticos}

\begin{lstlisting}
{
  "diagnostics": {
    "logLevel":          "off",
    "persistentLogs":    true,
    "interactionHistory": true,
    "email": "your-email@example.com"
  }
}
\end{lstlisting}

Valid log levels: \code{off}, \code{critical}, \code{warning},
\code{info}, \code{debug}, \code{full}.

% ── A.13  Troubleshooting installation ───────────────────────
\section{Problemas de instalación de solución de problemas}
\label{app:troubleshooting-install}

\subsection{FairWindSK no comienza después de la instalación}

\begin{enumerate}
  \item Run the binary from a terminal to see error output:
        \code{FairWindSK} (Linux/macOS) or
        \code{build\textbackslash FairWindSK.exe} (Windows).
  \item Check that Qt WebEngine libraries are available.
        On Linux, run \code{ldd build/FairWindSK | grep "not found"}.
  \item On Raspberry~Pi, check for the virtual keyboard CMake workaround
        (Section~\ref{app:rpi-station}).
\end{enumerate}

\subsection{No aparecen aplicaciones en el lanzador}

\begin{enumerate}
  \item Verify the Signal~K server address in
        \ui{Settings~> Connection}.
  \item Confirm the server exposes
        \code{/signalk/v1/apps/list} or \code{/skServer/webapps}.
  \item Check the Signal~K App Store --- at least one web app must be
instalado en el servidor para que aparezca en FairWindSK.
  \item Enable debug logging: add \code{debug=true} to
        \code{fairwindsk.ini} and relaunch.
Los registros incluyen intentos de conexión y detalles de descubrimiento de aplicaciones.
\end{enumerate}

\subsection{Construir fallas con el módulo Qt desaparecido}

\begin{enumerate}
  \item Verify the Qt version: \code{qmake6 --version} or
        \code{cmake --find-package -DNAME=Qt6 -DCOMPILER\_ID=GNU \ldots}
  \item Check that all required packages are installed
        (Sections~\ref{app:linux} and~\ref{app:rpi-station}).
  \item If \code{CMAKE\_PREFIX\_PATH} is needed, pass it explicitly:
        \code{cmake -S . -B build -DCMAKE\_PREFIX\_PATH=/path/to/Qt/6.x.x/...}
\end{enumerate}

\subsection{FairWindSK no arranca en Raspberry Pi}

\begin{enumerate}
  \item Verify the autostart file was installed:
\begin{lstlisting}
ls /etc/xdg/autostart/fairwindsk-startup.desktop
\end{lstlisting}
  \item Verify the \code{fairwindsk-startup} script is executable
y en PATH:
\begin{lstlisting}
which fairwindsk-startup
\end{lstlisting}
  \item Run the startup helper manually from a terminal to check for errors:
\begin{lstlisting}
fairwindsk-startup
\end{lstlisting}
  \item Ensure the desktop session supports XDG autostart
(el escritorio Raspberry~Pi~OS predeterminado lo hace).
\end{enumerate}
